\SourcePage{603}{606}
\section{Hyperfine structure of molecular levels}

The hyperfine structure of molecular energy levels has the same physical
origin as the hyperfine structure of atomic levels.

For the overwhelming majority of molecules the total electron spin is zero.
Their principal source of hyperfine level splitting is the quadrupole
interaction of the nuclei with the electrons. Naturally, only nuclei whose
spin $i$ is neither 0 nor $1/2$ participate in this interaction; otherwise the
quadrupole moment vanishes.

Because nuclear motion in a molecule is comparatively slow, the molecular
average of the quadrupole-interaction operator is taken in two stages. First
it is averaged over the electronic state with the nuclei held fixed, and then
over the rotation of the molecule.

Consider first a diatomic molecule. After the first averaging, the interaction
of each nucleus with the electrons is represented by an operator proportional
to the scalar $\hat Q_{ik}n_in_k$. This scalar is formed from the operator of
the nuclear quadrupole-moment tensor and the unit vector $\mathbf n$ along the
molecular axis, the sole quantity that specifies the orientation of the
molecule relative to the direction of the nuclear spin. Since
$\hat Q_{ii}=0$, the operator may be written
\begin{equation}
 b\hat i_i\hat i_k\left(n_in_k-\frac13\delta_{ik}\right);
 \tag{122.1}
\end{equation}
for a prescribed value of the projection $i_\zeta$ of the nuclear spin on the
molecular axis, this quantity equals
$b[i_\zeta^2-\frac13i(i+1)]$.

On averaging (122.1) over the molecular rotation, however, it is expressed in
terms of the operator $\hat{\mathbf K}$ of the conserved rotational angular
momentum. The product $n_in_k$ is averaged by the formula obtained in the
problem to \S~29 (with $\mathbf K$ in place of $\mathbf l$), giving
\begin{equation}
 -\frac{b}{(2K-1)(2K+3)}\hat i_i\hat i_k
 \left[\hat K_i\hat K_k+\hat K_k\hat K_i
 -\frac23\delta_{ik}K(K+1)\right].
 \tag{122.2}
\end{equation}
The eigenvalues of this operator are found in the same way as was described
for the operator (121.6).

For a polyatomic molecule, (122.1) is replaced in general by an operator of
the form
\begin{equation}
 b_{ik}\hat i_i\hat i_k,
 \tag{122.3}
\end{equation}

\SourcePage{604}{607}
where $b_{ik}$ is a traceless tensor constituting a particular characteristic
of the molecular electronic state. After averaging over the molecular
rotation, this tensor is expressed in terms of the molecule's total rotational
angular momentum $\mathbf J$ by a formula of the form
\begin{equation}
 \overline{b_{ik}}=b\left[\hat J_i\hat J_k+\hat J_k\hat J_i
 -\frac23J(J+1)\delta_{ik}\right].
 \tag{122.4}
\end{equation}

In principle the coefficient $b$ can be expressed through the components of
the tensor $b_{ik}$ referred to the principal axes of inertia $\xi$, $\eta$,
$\zeta$ of the molecule. Since these axes are rigidly fixed in the molecule,
the components $b_{\xi\xi},\ldots$ are molecular characteristics unaffected
by the averaging. To do this, consider the scalar
$\overline{b_{ik}J_iJ_k}$. Calculation by means of (122.4) gives
\begin{equation}
 \overline{b_{ik}J_iJ_k}
 =bJ(J+1)\left[\frac43J(J+1)-1\right]
 \tag{122.5}
\end{equation}
(the calculation is analogous to that in the problem to \S~29). On the other
hand, expanding the tensor product in the $\xi$, $\eta$, $\zeta$ axes gives
\begin{equation}
 \overline{b_{ik}J_iJ_k}
 =b_{\xi\xi}\overline{J_\xi^2}
 +b_{\eta\eta}\overline{J_\eta^2}
 +b_{\zeta\zeta}\overline{J_\zeta^2}.
 \tag{122.6}
\end{equation}
Here it has been used that the mean values of the products
$J_\xi J_\zeta,\ldots$ vanish.\footnote{Indeed, in a representation in which
the matrix of one component of $\mathbf J$ (say $J_\zeta$) is diagonal, the
matrices of the products $J_\xi J_\zeta$, $J_\eta J_\zeta$ have elements only
when the quantum number $k$ changes by 1; the wave functions of stationary
states of an asymmetric top, however, contain functions $\psi_{Jk}$ with
values of $k$ differing by an even number (see \S~103).} The mean values of
the squares $J_\xi^2,\ldots$ can in principle be calculated from the wave
functions of the corresponding rotational states of the top. In particular,
for a symmetric top one simply has
\[
 \overline{J_\zeta^2}=k^2,\qquad
 \overline{J_\xi^2}=\overline{J_\eta^2}
 =\frac12[J(J+1)-k^2].
\]

When the nuclear spins are $1/2$, there is no quadrupole interaction. One of
the principal sources of hyperfine splitting in this case is the direct
magnetic interaction between the nuclear magnetic moments. The interaction
operator of two magnetic moments
$\boldsymbol\mu_1=\mu_1\hat{\mathbf i}_1/i_1$,
$\boldsymbol\mu_2=\mu_2\hat{\mathbf i}_2/i_2$ is
\[
 \frac{\mu_1\mu_2}{i_1i_2r^3}
 [\hat{\mathbf i}_1\hat{\mathbf i}_2
 -3(\hat{\mathbf i}_1\mathbf n)(\hat{\mathbf i}_2\mathbf n)].
\]

\SourcePage{605}{608}
To calculate the splitting energy, this operator must be averaged over the
molecular state in a manner similar to that described above.

When heavy atoms are present in the molecule, an indirect interaction of the
nuclear moments mediated by the electron shell makes a contribution to the
hyperfine splitting comparable with that of the direct interaction. Formally,
this interaction is a second-order perturbation-theory effect with respect to
the interaction of the nuclear spin with the electrons. From the results of
\S~121 it is readily found that the ratio of this effect to that of the direct
interaction of the nuclear moments is of order $(Ze^2/\hbar c)^2$; for large
$Z$ it is comparable with unity.

Finally, the interaction of a nuclear moment with molecular rotation also
makes a certain contribution to the hyperfine splitting of molecular levels.
A rotating molecule, being a moving system of charges, produces a magnetic
field. This field can be calculated by the familiar electrodynamic formulas
from the prescribed current density
$\mathbf j=\rho(\boldsymbol\Omega\times\mathbf r)$, where $\rho$ is the
density of the charges (electrons and nuclei) in the molecule at rest and
$\boldsymbol\Omega$ is its angular velocity. The level splitting is obtained
as the energy of the nuclear magnetic moment in this field, with the
components of the molecular angular velocity expressed through the components
of its angular momentum (cf. \S~103).
